\section{Homological Evaluation, Visible Quotients, and Pure-Ext Blindness}\label{sec:homological-visibility}

To isolate the part of a branch obstruction that belongs to the cohomology class itself, we make the homological evaluation map primary. The intrinsic quotient will be defined from this map and then proved invariant under the two controlled changes allowed in this paper. Homologically equivalent finite presentations are handled by the presentation-invariance theorem. Coefficient retivializations on common refinements are handled by the transport-invariance theorem. Only after these invariance results do we recover the quotient through class-admissible character channels.

Here ``intrinsic'' means intrinsic to the chosen branch obstruction class $\omega$ on its allowed finite-presentation transport class: homological finite-presentation changes with $H_2$-isomorphic nerves, together with coefficient retivializations by finite-group isomorphisms on common refinements. It does not mean intrinsic to the underlying set-valued local object functor before such a realization-theoretic choice has been made.

\begin{definition}[Finite nerve presentation of the obstruction class]\label{def:finite-nerve-presentation}
Work in the setting of \Cref{def:cech-gluing-obstruction} and fix a visible
branch
\[
v\in \Vis_{p,s}(r).
\]
A \emph{finite nerve presentation} of the branch obstruction class of $v$
consists of the following data:
\begin{enumerate}[label=(\roman*),nosep]
\item a finite gerbe-adapted cover
\[
\mathcal U=\{U_i\to a\}_{i\in I}\in \mathcal A_v;
\]
\item a finite abelian group $A$ together with a trivialization of the branch
band on every nonempty simplex of the \v{C}ech nerve of $\mathcal U$, so that
the corresponding simplicial \v{C}ech complex has constant coefficients in $A$;
\item a class
\[
\omega\in H^2(N(\mathcal U),A)
\]
representing the restriction of the filtered branch obstruction class
$\omega_v$ to the cover $\mathcal U$ under the identification of
\Cref{thm:cech-bridge-compatible-realizations}.
\end{enumerate}
For such a presentation, the universal coefficient theorem for the finite
simplicial set $N(\mathcal U)$ gives a natural short exact sequence
\[
0\longrightarrow
\operatorname{Ext}^1(H_1(N(\mathcal U),\mathbb Z),A)
\longrightarrow
H^2(N(\mathcal U),A)
\xrightarrow{\ \mathrm{ev}\ }
\operatorname{Hom}(H_2(N(\mathcal U),\mathbb Z),A)
\longrightarrow 0
\]
\cite[Sec.~3A]{Hatcher2002}, \cite[Sec.~3.5]{Weibel1994}.
The image of $\omega$ under the displayed map is denoted
\[
\mathrm{ev}_{\omega}:H_2(N(\mathcal U),\mathbb Z)\to A
\]
and called the \emph{homological evaluation map} of the finite nerve
presentation.
\end{definition}

\begin{definition}[Intrinsic hidden subgroup and visible quotient]\label{def:intrinsic-visible-quotient}
Let $(\mathcal U,A,\omega)$ be a finite nerve presentation. The
\emph{intrinsic hidden subgroup} of $\omega$ is the image subgroup
\[
K_{\omega}:=\operatorname{Im}(\mathrm{ev}_{\omega})\le A.
\]
The associated \emph{intrinsic visible quotient} is the quotient finite
abelian group
\[
A_{\mathrm{vis}}^{\omega}:=A/K_{\omega}.
\]
We write
\[
q_{\omega}:A\twoheadrightarrow A_{\mathrm{vis}}^{\omega}
\]
for the quotient map. Thus the part of the obstruction class detected on
$H_2(N(\mathcal U),\mathbb Z)$ is recorded by the subgroup $K_{\omega}$, and
the coefficient group visible after modding out this detected part is
$A_{\mathrm{vis}}^{\omega}$.
\end{definition}

\begin{proposition}[Canonical quotient sequence of the homological evaluation]\label{prop:visible-quotient-quotient-sequence}
Let $(\mathcal U,A,\omega)$ be a finite nerve presentation. Then there is a
canonical short exact sequence of finite abelian groups
\[
0\longrightarrow K_{\omega}\longrightarrow A
\xrightarrow{\ q_{\omega}\ } A_{\mathrm{vis}}^{\omega}\longrightarrow 0,
\]
with
\[
K_{\omega}=\operatorname{Im}(\mathrm{ev}_{\omega}).
\]
Equivalently,
\[
q_{\omega}\circ \mathrm{ev}_{\omega}=0,
\]
and for every homomorphism $\pi:A\to Q$ into an abelian group $Q$, the
following are equivalent:
\begin{enumerate}[label=(\roman*),nosep]
\item $\pi\circ \mathrm{ev}_{\omega}=0$;
\item $K_{\omega}\subseteq \ker(\pi)$;
\item there is a unique homomorphism $\overline\pi:A_{\mathrm{vis}}^{\omega}\to Q$
such that
\[
\pi=\overline\pi\circ q_{\omega}.
\]
\end{enumerate}
\end{proposition}

\begin{proof}
Since $K_{\omega}$ is defined to be the image of $\mathrm{ev}_{\omega}$, the
sequence
\[
0\longrightarrow K_{\omega}\longrightarrow A\longrightarrow A/K_{\omega}
\longrightarrow 0
\]
is the standard quotient exact sequence, and $A/K_{\omega}$ is precisely
$A_{\mathrm{vis}}^{\omega}$ by \Cref{def:intrinsic-visible-quotient}. In
particular, the quotient map $q_{\omega}$ kills $K_{\omega}$, hence kills
$\operatorname{Im}(\mathrm{ev}_{\omega})$, so
$q_{\omega}\circ \mathrm{ev}_{\omega}=0$.

For a homomorphism $\pi:A\to Q$, the condition
$\pi\circ \mathrm{ev}_{\omega}=0$ means exactly that $\pi$ vanishes on the
image of $\mathrm{ev}_{\omega}$, namely on $K_{\omega}$. This proves
\textup{(i)}$\Leftrightarrow$\textup{(ii)}. The equivalence with
\textup{(iii)} is the universal property of the quotient group
$A\twoheadrightarrow A/K_{\omega}=A_{\mathrm{vis}}^{\omega}$.
\end{proof}

\begin{proposition}[Coefficient transport of the hidden subgroup and visible quotient]\label{prop:visible-quotient-coefficient-transport}
Let $(\mathcal U,A,\omega)$ be a finite nerve presentation, let
\[
\phi:A\xrightarrow{\sim} A'
\]
be an isomorphism of finite abelian groups, and write
\[
\omega':=\phi_*(\omega)\in H^2(N(\mathcal U),A').
\]
Then
\[
\mathrm{ev}_{\omega'}=\phi\circ \mathrm{ev}_{\omega},
\qquad
K_{\omega'}=\phi(K_{\omega}),
\]
and $\phi$ induces an isomorphism
\[
\overline\phi:A_{\mathrm{vis}}^{\omega}\xrightarrow{\sim}A_{\mathrm{vis}}^{\omega'}
\]
such that
\[
q_{\omega'}\circ \phi=\overline\phi\circ q_{\omega}.
\]
\end{proposition}

\begin{proof}
Naturality of the universal coefficient sequence with respect to the
coefficient homomorphism $\phi$ gives
\[
\mathrm{ev}_{\omega'}=\mathrm{ev}_{\phi_*(\omega)}=\phi\circ \mathrm{ev}_{\omega}.
\]
Taking images yields
\[
K_{\omega'}=\operatorname{Im}(\mathrm{ev}_{\omega'})
=\operatorname{Im}(\phi\circ \mathrm{ev}_{\omega})
=\phi(\operatorname{Im}(\mathrm{ev}_{\omega}))
=\phi(K_{\omega}).
\]
Since $\phi$ carries $K_{\omega}$ isomorphically onto $K_{\omega'}$, it
induces an isomorphism of quotients
\[
\overline\phi:A/K_{\omega}\xrightarrow{\sim}A'/K_{\omega'}.
\]
By \Cref{def:intrinsic-visible-quotient}, these quotients are exactly
$A_{\mathrm{vis}}^{\omega}$ and $A_{\mathrm{vis}}^{\omega'}$. The displayed
intertwining relation is the standard compatibility of an induced quotient map
with the original homomorphism $\phi$.
\end{proof}

\begin{definition}[Homologically equivalent finite nerve presentations]\label{def:homologically-equivalent-finite-presentations}
Let
\[
(\mathcal U,A,\omega)
\qquad\text{and}\qquad
(\mathcal U',A,\omega')
\]
be two finite nerve presentations of the same branch obstruction, with the same
finite coefficient group $A$ coming from the chosen stack-level band
trivialization. They are \emph{homologically equivalent} if there is a third
finite nerve presentation
\[
(\mathcal W,A,\widetilde\omega)
\]
which refines both and whose refinement maps of nerves
\[
\rho:N(\mathcal W)\longrightarrow N(\mathcal U),
\qquad
\rho':N(\mathcal W)\longrightarrow N(\mathcal U')
\]
satisfy
\[
\rho^*(\omega)=\widetilde\omega=(\rho')^*(\omega')
\]
in $H^2(N(\mathcal W),A)$ and induce isomorphisms
\[
H_2(\rho):H_2(N(\mathcal W),\mathbb Z)\xrightarrow{\sim}
H_2(N(\mathcal U),\mathbb Z),
\]
\[
H_2(\rho'):H_2(N(\mathcal W),\mathbb Z)\xrightarrow{\sim}
H_2(N(\mathcal U'),\mathbb Z).
\]
A direct finite refinement is an allowed presentation change when its nerve map
satisfies the same $H_2$-isomorphism condition and pulls the obstruction class
to the refined obstruction class.
\end{definition}

\begin{theorem}[Presentation invariance of the visible quotient]\label{thm:visible-quotient-presentation-invariance}
Homologically equivalent finite nerve presentations have the same hidden
subgroup and the same visible quotient:
\[
K_{\omega}=K_{\omega'}
\qquad\text{and}\qquad
A_{\mathrm{vis}}^{\omega}=A_{\mathrm{vis}}^{\omega'}
\]
as subgroups and quotients of the common coefficient group $A$.
\end{theorem}

\begin{proof}
Let $(\mathcal W,A,\widetilde\omega)$ be a common homological refinement as in
\Cref{def:homologically-equivalent-finite-presentations}. Naturality of the
universal coefficient sequence for the simplicial maps $\rho$ and $\rho'$ gives
\[
\mathrm{ev}_{\widetilde\omega}
=
\mathrm{ev}_{\omega}\circ H_2(\rho)
=
\mathrm{ev}_{\omega'}\circ H_2(\rho').
\]
Because $H_2(\rho)$ and $H_2(\rho')$ are isomorphisms, composing with either of
them does not change the image subgroup in $A$. Hence
\[
\mathrm{Im}(\mathrm{ev}_{\omega})
=
\mathrm{Im}(\mathrm{ev}_{\widetilde\omega})
=
\mathrm{Im}(\mathrm{ev}_{\omega'}).
\]
By \Cref{def:intrinsic-visible-quotient}, this says
$K_{\omega}=K_{\omega'}$. Taking the quotient of the same finite abelian group
$A$ by this common subgroup gives
$A_{\mathrm{vis}}^{\omega}=A_{\mathrm{vis}}^{\omega'}$.
\end{proof}

\begin{definition}[Transport-equivalent finite nerve presentations]\label{def:transport-equivalent-finite-presentations}
Let
\[
(\mathcal U,A,\omega)
\qquad\text{and}\qquad
(\mathcal U',A',\omega')
\]
be finite nerve presentations of the same branch obstruction. They are
\emph{transport-equivalent} if there is a third finite nerve presentation
\[
(\mathcal W,B,\widetilde\omega)
\]
with refinement maps of nerves
\[
\rho:N(\mathcal W)\longrightarrow N(\mathcal U),
\qquad
\rho':N(\mathcal W)\longrightarrow N(\mathcal U')
\]
and coefficient isomorphisms
\[
\alpha:B\xrightarrow{\sim} A,
\qquad
\alpha':B\xrightarrow{\sim} A'
\]
such that
\[
\alpha_*(\widetilde\omega)=\rho^*(\omega),
\qquad
\alpha'_*(\widetilde\omega)=(\rho')^*(\omega')
\]
and
\[
H_2(\rho):H_2(N(\mathcal W),\mathbb Z)\xrightarrow{\sim}
H_2(N(\mathcal U),\mathbb Z),
\]
\[
H_2(\rho'):H_2(N(\mathcal W),\mathbb Z)\xrightarrow{\sim}
H_2(N(\mathcal U'),\mathbb Z).
\]
This is the descent relation produced by passing to a common finite refinement
and retivializing the constant coefficient group by finite-group isomorphisms.
\end{definition}

\begin{theorem}[Visible quotient invariance under coefficient transport]\label{thm:visible-quotient-transport-invariance}
Transport-equivalent finite nerve presentations have the same hidden subgroup
and the same visible quotient up to the induced coefficient transport
isomorphism. More precisely, in the setting of
\Cref{def:transport-equivalent-finite-presentations}, the isomorphism
\[
\psi:=\alpha'\circ \alpha^{-1}:A\xrightarrow{\sim} A'
\]
satisfies
\[
\psi(K_{\omega})=K_{\omega'},
\]
and induces an isomorphism
\[
\overline\psi:A_{\mathrm{vis}}^{\omega}\xrightarrow{\sim}
A_{\mathrm{vis}}^{\omega'}.
\]
\end{theorem}

\begin{proof}
By naturality with respect to the simplicial map $\rho$ and then with respect
to the coefficient isomorphism $\alpha$,
\[
\alpha\circ \mathrm{ev}_{\widetilde\omega}
=
\mathrm{ev}_{\alpha_*(\widetilde\omega)}
=
\mathrm{ev}_{\rho^*(\omega)}
=
\mathrm{ev}_{\omega}\circ H_2(\rho).
\]
Since $H_2(\rho)$ is an isomorphism, taking images gives
\[
\alpha(K_{\widetilde\omega})=K_{\omega}.
\]
The same argument with $\rho'$ and $\alpha'$ yields
\[
\alpha'(K_{\widetilde\omega})=K_{\omega'}.
\]
Therefore
\[
\psi(K_{\omega})
=
\alpha'\alpha^{-1}(K_{\omega})
=
\alpha'(K_{\widetilde\omega})
=
K_{\omega'}.
\]
Because $\psi$ carries $K_{\omega}$ isomorphically onto $K_{\omega'}$, it
induces an isomorphism of quotient groups
\[
A/K_{\omega}\xrightarrow{\sim} A'/K_{\omega'},
\]
which is exactly the claimed isomorphism
\[
\overline\psi:A_{\mathrm{vis}}^{\omega}\xrightarrow{\sim}
A_{\mathrm{vis}}^{\omega'}
\]
after invoking \Cref{def:intrinsic-visible-quotient}.
\end{proof}

\begin{remark}[No naive refinement invariance]\label{rem:no-naive-visible-quotient-invariance}
\Cref{thm:visible-quotient-presentation-invariance,thm:visible-quotient-transport-invariance}
are the invariance statements used in this paper. They apply to finite
refinements that preserve the $H_2$-detected part of the obstruction through
the displayed homology isomorphisms, together with coefficient retivializations
by finite-group isomorphisms on a common refinement. An arbitrary replacement
of the finite nerve may change $H_2$ and therefore may change the image of the
homological evaluation map; likewise an arbitrary change of coefficients that is
not related by such an isomorphism is not treated as an invariance of
$A_{\mathrm{vis}}^{\omega}$.
\end{remark}

\begin{definition}[Class-admissible characters]\label{def:class-admissible-character}
In the setting of \Cref{def:finite-nerve-presentation}, write
\[
\widehat A:=\mathrm{Hom}(A,\mathbb T)
\]
for the character group of the finite abelian coefficient group $A$, and let
\[
X_{\omega}:=\{\chi\in\widehat A\mid \chi_*(\omega)=0\}.
\]
Elements of $X_{\omega}$ are called \emph{class-admissible} for $\omega$. Thus
a character is class-admissible exactly when pushing the obstruction class
forward along that character kills it in $H^2(N(\mathcal U),\mathbb T)$.
\end{definition}

\begin{theorem}[Character-theoretic detection of the intrinsic visible quotient]\label{thm:intrinsic-visible-quotient}
Let $(\mathcal U,A,\omega)$ be a finite nerve presentation and let
$q_{\omega}:A\to A_{\mathrm{vis}}^{\omega}$ be the quotient map of
\Cref{def:intrinsic-visible-quotient}. For a character
\[
\chi\in \widehat A=\operatorname{Hom}(A,\mathbb T),
\]
the following are equivalent:
\begin{enumerate}[label=(\roman*),nosep]
\item $\chi$ is class-admissible for $\omega$;
\item $\chi\circ \mathrm{ev}_{\omega}=0$;
\item $\chi$ is trivial on $K_{\omega}$;
\item there is a unique character
\[
\overline\chi:A_{\mathrm{vis}}^{\omega}\to\mathbb T
\]
such that
\[
\chi=\overline\chi\circ q_{\omega}.
\]
\end{enumerate}
Equivalently,
\[
X_{\omega}=q_{\omega}^*(\widehat{A_{\mathrm{vis}}^{\omega}})
\subseteq \widehat A.
\]
\end{theorem}

\begin{proof}
Because $\mathbb T$ is a divisible abelian group, it is an injective
$\mathbb Z$-module, so
\[
\operatorname{Ext}^1(H_1(N(\mathcal U),\mathbb Z),\mathbb T)=0
\]
\cite[Cor.~2.3.15]{Weibel1994}. Applying the universal coefficient theorem to
the coefficient group $\mathbb T$ therefore gives an isomorphism
\[
H^2(N(\mathcal U),\mathbb T)
\xrightarrow{\sim}
\operatorname{Hom}(H_2(N(\mathcal U),\mathbb Z),\mathbb T).
\]
By naturality of the universal coefficient sequence with respect to the
coefficient homomorphism $\chi:A\to\mathbb T$, the class $\chi_*(\omega)$ maps
to the composite
\[
\chi\circ \mathrm{ev}_{\omega}:H_2(N(\mathcal U),\mathbb Z)\to \mathbb T.
\]
Hence $\chi_*(\omega)=0$ if and only if
$\chi\circ \mathrm{ev}_{\omega}=0$, which proves
\textup{(i)}$\Leftrightarrow$\textup{(ii)}.

Since $K_{\omega}=\operatorname{Im}(\mathrm{ev}_{\omega})$,
\textup{(ii)} holds exactly when $\chi$ vanishes on the image of
$\mathrm{ev}_{\omega}$, which is \textup{(iii)}. A character of $A$ is trivial
on $K_{\omega}$ if and only if it factors uniquely through the quotient
$A\to A/K_{\omega}=A_{\mathrm{vis}}^{\omega}$, which is \textup{(iv)}.
The displayed equality of character groups is a restatement of the same
factorization criterion.
\end{proof}

\begin{corollary}[Character detection of the intrinsic visible quotient]\label{cor:intrinsic-character-detection}
Let $(\mathcal U,A,\omega)$ be a finite nerve presentation. Then the
class-admissible characters detect exactly the quotient
$A_{\mathrm{vis}}^{\omega}$:
\begin{enumerate}[label=(\roman*),nosep]
\item if two elements of $A$ have the same image in $A_{\mathrm{vis}}^{\omega}$,
then every class-admissible character takes the same value on them;
\item if
\[
\bar a\neq \bar b\in A_{\mathrm{vis}}^{\omega},
\]
then there is a class-admissible character
\[
\chi\in X_{\omega}
\]
whose factor
\[
\overline\chi:A_{\mathrm{vis}}^{\omega}\to\mathbb T
\]
satisfies $\overline\chi(\bar a)\neq \overline\chi(\bar b)$.
\end{enumerate}
Equivalently, $A_{\mathrm{vis}}^{\omega}$ is the largest quotient of $A$ whose
points are separated by the class-admissible character channels.
\end{corollary}

\begin{proof}
If $a,b\in A$ have the same image in $A_{\mathrm{vis}}^{\omega}$, then every
class-admissible character factors through $q_{\omega}$ by
\Cref{thm:intrinsic-visible-quotient}, so it takes the same value on $a$ and
$b$. This proves \textup{(i)}.

For \textup{(ii)}, the difference $\bar a-\bar b$ is a nonzero element of the
finite abelian group $A_{\mathrm{vis}}^{\omega}$. Finite abelian groups are
separated by characters \cite[Chap.~2]{Terras1999}, so there exists a
character
\[
\overline\chi:A_{\mathrm{vis}}^{\omega}\to\mathbb T
\]
with $\overline\chi(\bar a)\neq \overline\chi(\bar b)$. Then
$\chi:=\overline\chi\circ q_{\omega}$ is class-admissible by
\Cref{thm:intrinsic-visible-quotient}. The final sentence is exactly the
combination of \textup{(i)} and \textup{(ii)}.
\end{proof}

\begin{theorem}[Initiality among pure-extension quotients]\label{thm:intrinsic-pure-ext-initiality}
Let $(\mathcal U,A,\omega)$ be a finite nerve presentation and let
$q_{\omega}:A\to A_{\mathrm{vis}}^{\omega}$ be the intrinsic visible quotient.
For a quotient homomorphism
\[
\pi:A\twoheadrightarrow Q
\]
onto a finite abelian group $Q$, the following are equivalent:
\begin{enumerate}[label=(\roman*),nosep]
\item the pushed-forward class $\pi_*(\omega)$ lies in the pure
$\operatorname{Ext}$-part of the universal coefficient sequence for
$H^2(N(\mathcal U),Q)$;
\item
\[
\mathrm{ev}_{\pi_*(\omega)}=0;
\]
\item
\[
\pi\circ \mathrm{ev}_{\omega}=0;
\]
\item
\[
K_{\omega}\subseteq \ker(\pi).
\]
\end{enumerate}
When these conditions hold, there is a unique quotient map
\[
\overline\pi:A_{\mathrm{vis}}^{\omega}\twoheadrightarrow Q
\]
such that
\[
\pi=\overline\pi\circ q_{\omega}.
\]
Thus $A_{\mathrm{vis}}^{\omega}$ is initial among coefficient quotients through
which the obstruction becomes purely extension-theoretic.
\end{theorem}

\begin{proof}
For the coefficient map $\pi:A\to Q$, naturality of the universal coefficient
sequence gives
\[
\mathrm{ev}_{\pi_*(\omega)}=\pi\circ \mathrm{ev}_{\omega}.
\]
Hence \textup{(ii)} and \textup{(iii)} are equivalent. Exactness of the
universal coefficient sequence for $Q$ identifies
\[
\ker\!\bigl(H^2(N(\mathcal U),Q)\to
\operatorname{Hom}(H_2(N(\mathcal U),\mathbb Z),Q)\bigr)
\]
with the image of
$\operatorname{Ext}^1(H_1(N(\mathcal U),\mathbb Z),Q)$, so
\textup{(i)}$\Leftrightarrow$\textup{(ii)}. Since
$K_{\omega}=\operatorname{Im}(\mathrm{ev}_{\omega})$,
\textup{(iii)} holds exactly when $\pi$ kills that image, which is
\textup{(iv)}. The unique factorization through $A/K_{\omega}$ is then the
universal property of the quotient.
\end{proof}

\begin{corollary}[Class-trivializing initiality when $H_1$ vanishes]\label{cor:h1-vanishing-class-trivializing}
Let $(\mathcal U,A,\omega)$ be a finite nerve presentation and assume
\[
H_1(N(\mathcal U),\mathbb Z)=0.
\]
Then, for every quotient homomorphism $\pi:A\twoheadrightarrow Q$, the
following are equivalent:
\begin{enumerate}[label=(\roman*),nosep]
\item $\pi_*(\omega)=0$ in $H^2(N(\mathcal U),Q)$;
\item $\pi$ factors uniquely through the intrinsic visible quotient
$q_{\omega}:A\to A_{\mathrm{vis}}^{\omega}$.
\end{enumerate}
In particular, when $H_1$ vanishes the intrinsic visible quotient is initial
among coefficient quotients that trivialize the class.
\end{corollary}

\begin{proof}
If $H_1(N(\mathcal U),\mathbb Z)=0$, then
\[
\operatorname{Ext}^1(H_1(N(\mathcal U),\mathbb Z),Q)=0
\]
for every $Q$. Hence in the universal coefficient sequence for $Q$ a class is
pure-$\operatorname{Ext}$ if and only if it is zero. The claim is therefore
exactly \Cref{thm:intrinsic-pure-ext-initiality}.
\end{proof}

\begin{remark}[Canonical visibility without a canonical splitting]\label{rem:no-canonical-splitting}
For a chosen finite nerve presentation, the quotient
$A_{\mathrm{vis}}^{\omega}$ and the subgroup
$K_{\omega}=\mathrm{Im}(\mathrm{ev}_{\omega})$ are canonical, and
\Cref{thm:visible-quotient-presentation-invariance,thm:visible-quotient-transport-invariance}
show that they descend to the allowed coefficient-transport class of the branch
obstruction. By contrast, the universal coefficient sequence need not split
canonically. Thus the theory isolates a canonical visible quotient and a
canonical blind subgroup on that class without requiring any noncanonical
decomposition of $\omega$ into visible and invisible summands.
\end{remark}

\begin{theorem}[Character-blind pure-extension obstructions]\label{thm:character-blind-obstructions}
Let $(\mathcal U,A,\omega)$ be a finite nerve presentation. The following are
equivalent:
\begin{enumerate}[label=(\roman*),nosep]
\item every character of $A$ is class-admissible for $\omega$;
\item
\[
X_{\omega}=\widehat A;
\]
\item
\[
K_{\omega}=0;
\]
\item
\[
\mathrm{ev}_{\omega}=0;
\]
\item $\omega$ lies in the image of
\[
\operatorname{Ext}^1(H_1(N(\mathcal U),\mathbb Z),A)\longrightarrow
H^2(N(\mathcal U),A).
\]
\end{enumerate}
When these conditions hold,
\[
A_{\mathrm{vis}}^{\omega}=A,
\]
so the obstruction is nonzero but completely invisible to every character
channel exactly when it is a pure $\operatorname{Ext}$-class.
\end{theorem}

\begin{proof}
The equivalence \textup{(i)}$\Leftrightarrow$\textup{(ii)} is only notation.
By \Cref{thm:intrinsic-visible-quotient}, every character is
class-admissible if and only if every character is trivial on $K_{\omega}$.
Characters separate points of finite abelian groups
\cite[Chap.~2]{Terras1999}, so this happens if and only if $K_{\omega}=0$.
Thus \textup{(ii)}$\Leftrightarrow$\textup{(iii)}.

Since $K_{\omega}=\operatorname{Im}(\mathrm{ev}_{\omega})$,
\textup{(iii)} and \textup{(iv)} are equivalent. Exactness of the universal
coefficient sequence identifies the kernel of the evaluation map with the image
of the $\operatorname{Ext}^1$-term, so \textup{(iv)}$\Leftrightarrow$\textup{(v)}.
If these conditions hold, then
$A_{\mathrm{vis}}^{\omega}=A/K_{\omega}=A$ by definition.
\end{proof}

\begin{corollary}[Automatic blindness when $H_2$ vanishes]\label{cor:h2-vanishing-blindness}
Let $(\mathcal U,A,\omega)$ be a finite nerve presentation. If
\[
H_2(N(\mathcal U),\mathbb Z)=0,
\]
then $\omega$ is completely character-blind. Equivalently,
\[
X_{\omega}=\widehat A,
\qquad
A_{\mathrm{vis}}^{\omega}=A.
\]
\end{corollary}

\begin{proof}
If $H_2(N(\mathcal U),\mathbb Z)=0$, then the homological evaluation map
\[
\mathrm{ev}_{\omega}:H_2(N(\mathcal U),\mathbb Z)\to A
\]
is necessarily zero. Apply \Cref{thm:character-blind-obstructions}.
\end{proof}

\begin{corollary}[No character-blind obstructions when $H_1$ is free]\label{cor:h1-free-no-blindness}
Let $(\mathcal U,A,\omega)$ be a finite nerve presentation and assume that
\[
H_1(N(\mathcal U),\mathbb Z)
\]
is free abelian. Then the following are equivalent:
\begin{enumerate}[label=(\roman*),nosep]
\item $\omega$ is completely character-blind;
\item $\omega=0$.
\end{enumerate}
In particular, no nonzero obstruction class is completely character-blind when
$H_1$ is free.
\end{corollary}

\begin{proof}
If $H_1(N(\mathcal U),\mathbb Z)$ is free, then
\[
\operatorname{Ext}^1(H_1(N(\mathcal U),\mathbb Z),A)=0.
\]
By \Cref{thm:character-blind-obstructions}, complete character-blindness is
equivalent to $\omega$ lying in this zero subgroup, hence to $\omega=0$.
\end{proof}

The following construction is stated before the two worked examples because it
is the point at which the conditional gerbe hypotheses become actual semantic
data. The order matters. We first choose an actual banded gerbe, then take its
presheaf of local isomorphism classes as the local object functor, and only then
read off the visible branch and its obstruction class. Thus the cocycles come
from gerbe data on the nerve rather than from arbitrarily declared comparison
elements on a bare set-valued presheaf, and the realization-prestack and
global-conservativity hypotheses used by the semantic theorems are verified
inside the construction.

\begin{lemma}[Good-cover site and nerve comparison tools]\label{lem:good-cover-site-tools}
Let $X$ be a compact triangulable space and let $\mathcal U$ be a finite good
cover. Then finite good covers refining $\mathcal U$ form a filtered cofinal
system after passing to open-star refinements of sufficiently fine barycentric
subdivisions. For the basis consisting of nonempty finite intersections of
members of those covers, restriction from the topological site of $X$ to the
basis site preserves sheaves and computes the same sheaf cohomology. Moreover,
for any finite good cover $\mathcal V$ and any constant finite abelian group
$A$, the \v{C}ech cochain complex of $A$ on $\mathcal V$ is the simplicial
cochain complex of $N(\mathcal V)$, and the nerve theorem identifies
$|N(\mathcal V)|$ with $X$ when $\mathcal V$ covers $X$.
\end{lemma}

\begin{proof}
Compact triangulable spaces admit finite triangulations, and open stars in
sufficiently fine barycentric subdivisions give finite good refinements of any
finite open cover; two finite good refinements are dominated by a common open
star refinement after a further subdivision
\cite[Sec.~2.1 and Chap.~4]{Hatcher2002}. This gives filteredness and
cofinality.

The nonempty finite intersections form a basis closed under finite
intersection after refinement. The comparison lemma for cohomology on a basis
says that sheaves and sheaf cohomology computed on this basis agree with those
on the corresponding topological subspace site
\cite[Chap.~III, Secs.~2--4]{MacLaneMoerdijk1994},
\cite[Sec.~4.1]{Bredon1997}. Since the members and nonempty finite
intersections of a good cover are contractible, the constant finite abelian
sheaf has vanishing positive cohomology on those pieces.

Finally, the good-cover hypothesis identifies the cover \v{C}ech complex with
the simplicial cochain complex of the nerve, and the nerve theorem gives the
homotopy equivalence between the realization of the nerve and the covered
space \cite[Sec.~4G]{Hatcher2002}.
\end{proof}

\begin{proposition}[Verified semantic realization of finite nerve obstructions]\label{prop:semantic-gerbe-realization}
Let $\mathcal U=\{a_i\to a\}_{i\in I}$ be a finite cover with \v{C}ech nerve $N(\mathcal U)$, let $A$ be a finite abelian group, and assume that the cover is gerbe-adapted for the constant band $A$:
\[
H^1(a_i,A)=0
\qquad\text{and}\qquad
H^1(a_{ij},A)=0
\]
for all $i,j$. Assume also that, in the slice site under consideration,
gerbe-adapted finite covers refining $\mathcal U$ form a cofinal family for the
\v{C}ech system used in \Cref{def:cech-gluing-obstruction}. Then every class
\[
\omega\in H^2(N(\mathcal U),A)
\]
is realized by a one-reference semantic instance with terminal slice object $a$
such that:
\begin{enumerate}[label=(\roman*),nosep]
\item the chosen stack-banded realization prestack is an explicitly chosen $A$-banded gerbe
$\mathfrak G_\omega$ with gerbe class $\omega$;
\item the local-object functor is
\[
F_{p,s}:=\pi_0^{\mathrm{pre}}(\mathfrak G_\omega);
\]
\item global conservativity at $a$ holds;
\item $\Vis_{p,s}(r)$ has one visible branch $v$, the component gerbe
$\mathfrak G_v$ is $\mathfrak G_\omega$, and the branch obstruction class is
exactly $\omega$;
\item if $\omega\neq 0$, then
\[
F_{p,s}(a)=\varnothing,
\qquad
\bigl(a_{p,s}F_{p,s}\bigr)(a)\neq\varnothing,
\qquad
M,p\Vdash \Null^{\mathrm{glue}}_s(r).
\]
\end{enumerate}
\end{proposition}

\begin{proof}
The cover is not merely a formal finite nerve: by hypothesis it belongs to the
cofinal class on which the branch obstruction is computed. Choose an
$A$-valued \v{C}ech $2$-cocycle $g=(g_{ijk})$ representing $\omega$. By the
standard classification of banded gerbes by degree-$2$ \v{C}ech classes
\cite[Chap.~IV]{Giraud1971}, \cite[Sec.~11]{StacksProject}, there is an
$A$-banded gerbe $\mathfrak G_\omega$ on the corresponding slice site whose
gerbe class is $\omega$, and it is neutral exactly when $\omega=0$.

Concretely, one may start with formal local objects $x_i$ over the cover
pieces, formal edge isomorphisms $\phi_{ij}:x_i|_{a_{ij}}\to x_j|_{a_{ij}}$,
and automorphism sheaf $A$, and impose on triple overlaps the relation
\[
\phi_{jk}\circ\phi_{ij}
=
\phi_{ik}\circ g_{ijk},
\]
using the sign convention of \Cref{thm:cech-bridge-compatible-realizations}. The cocycle
identity is exactly the associativity condition on quadruple overlaps, and
stackifying this descent presentation gives the same $A$-banded gerbe
$\mathfrak G_\omega$. Replacing $g$ by a cohomologous cocycle changes the edge
trivializations by the inverse-coboundary convention of
\Cref{lem:change-edge-trivializations}, so the resulting gerbe depends only on
$\omega$.

This construction gives an explicit semantic instance of the hypotheses used
below. The imposed $H^1$-vanishing is precisely the local condition used in
\Cref{thm:cech-bridge-compatible-realizations} to obtain edge isomorphisms and a
well-defined obstruction representative, and the cofinality assumption is the
cofinality requirement in \Cref{def:cech-gluing-obstruction}. Take a
one-reference situation whose terminal slice object is $a$, set the chosen
realization prestack to be the stack
\[
\mathfrak P_r:=\mathfrak G_\omega,
\]
and define the local object functor by
\[
F_{p,s}:=\pi_0^{\mathrm{pre}}(\mathfrak G_\omega).
\]
This is a legitimate local-object functor: $F_{p,s}(U)$ records the isomorphism classes of objects of the gerbe that actually exist over $U$. Since $\mathfrak G_\omega$ is already a stack in groupoids, $\mathfrak P_r$ is a realization prestack, its stackification morphism is the identity, and the stackification banding required in \Cref{def:realization-prestack}, item~\textup{(ii)} is just the given $A$-banding of $\mathfrak G_\omega$ on every object. Global conservativity at $a$ follows immediately. Thus the realization input is not inferred from the band; it is the displayed stack $\mathfrak G_\omega$, with its object sheaves, isomorphism sheaves, composition, and stack-level banding.

The semantic hypotheses used later are therefore verified in this construction rather than assumed silently:
\begin{enumerate}[label=(\alph*),nosep]
\item the realization prestack is the explicitly chosen stack $\mathfrak G_\omega$;
\item the band is the automorphism sheaf of that stack, identified with the constant sheaf $A$ on every object of the stackification;
\item global conservativity holds because stackification is the identity;
\item the gerbe-adapted cover hypotheses are exactly the displayed $H^1$-vanishing conditions, and their cofinality is part of the chosen slice-site structure.
\end{enumerate}

The sheaf of connected components $\pi_0(\mathfrak G_\omega)$ is the terminal sheaf, because a gerbe is locally nonempty and locally connected. Hence
\[
\bigl(a_{p,s}F_{p,s}\bigr)(a)\cong \pi_0(\mathfrak G_\omega)(a)\cong \{*\},
\]
so $\Vis_{p,s}(r)$ has one branch $v$. The component gerbe $\mathfrak G_v$ is
$\mathfrak G_\omega$, and the branch obstruction class is exactly $\omega$. If
$\omega\neq 0$, then $\mathfrak G_\omega(a)=\varnothing$ by non-neutrality, so
$F_{p,s}(a)=\varnothing$ while its sheafification has the displayed global
section. Thus the reference is compatibly locally sectionable but not globally
sectionable, and \Cref{thm:gerbe-null-semantics} gives a genuine gluing-level
absence for this semantic instance. The later $S^2$-type and
$\mathbb{R}P^2$-type examples are obtained by applying this construction to the
indicated finite good covers and cohomology classes.
\end{proof}

\begin{proposition}[Good-cover semantic slice]\label{prop:good-cover-semantic-slice}
Let $X$ be either $S^2$ or $\mathbb{R}P^2$, and let $\mathcal U$ be a finite
good cover of $X$. There is a small preorder site
\[
\mathcal C_X(\mathcal U)
\]
with terminal object $X$ such that:
\begin{enumerate}[label=(\roman*),nosep]
\item the objects are $X$ together with all nonempty finite intersections of
members of finite good covers of $X$ refining $\mathcal U$, and morphisms are
inclusions;
\item covers are finite families of inclusions which cover the target and
admit a finite good refinement by objects of this basis;
\item the finite good covers refining $\mathcal U$ form a filtered cofinal
family of covers of the terminal object $X$, and any finite family of such
covers has a finite common good refinement;
\item for every finite abelian group $A$, every finite good cover in this
cofinal family is gerbe-adapted for the constant band $A$:
\[
H^1(V_i,A)=0,
\qquad
H^1(V_i\cap V_j,A)=0
\]
for all cover pieces and pairwise overlaps;
\item for every finite good cover $\mathcal V$ in the cofinal family, the
\v{C}ech cochain complex of the constant band $A$ on $\mathcal V$ is the
simplicial cochain complex of the nerve $N(\mathcal V)$. In particular,
\[
\check H^2(\mathcal V,A)\cong H^2(N(\mathcal V),A),
\]
and for the original cover $\mathcal U$ the nerve theorem identifies
$H^*(N(\mathcal U),A)$ with $H^*(X,A)$.
\end{enumerate}
\end{proposition}

\begin{proof}
Both $S^2$ and $\mathbb{R}P^2$ are compact triangulable spaces. Starting from
the finite good cover $\mathcal U$, let $\mathcal G_X(\mathcal U)$ be the class
of finite good covers of $X$ refining $\mathcal U$, together with
$\mathcal U$ itself. Standard barycentric subdivision and open-star
constructions for compact triangulated spaces give finite good refinements of
any finite open cover and, given finitely many finite good covers, a further
finite good cover refining all of them by
\Cref{lem:good-cover-site-tools}. Hence
$\mathcal G_X(\mathcal U)$ is filtered under refinement and is cofinal among
the terminal covers considered below.

Define the object set of $\mathcal C_X(\mathcal U)$ to be $X$ and all nonempty
finite intersections of members of covers in $\mathcal G_X(\mathcal U)$. Put a
unique morphism $W\to W'$ exactly when $W\subseteq W'$. Thus $X$ is terminal.
A finite family $\{W_\ell\to W\}$ is declared to cover $W$ when the open sets
$W_\ell$ cover $W$ and the family admits a finite good refinement by objects of
the same basis. The identity cover is immediate. Pullback stability follows by
intersecting a good refinement with the smaller basis object and then replacing
the resulting finite cover by a finite good basis refinement. Transitivity is
obtained by composing the finite good refinements. Thus these families form a
Grothendieck pretopology on the preorder. By definition and the filteredness
proved above, the covers in $\mathcal G_X(\mathcal U)$ are cofinal among covers
of the terminal object, and finite common refinements remain in
$\mathcal G_X(\mathcal U)$.

If $\mathcal V=\{V_i\}$ is a finite good cover in
$\mathcal G_X(\mathcal U)$, then every $V_i$ and every nonempty
$V_i\cap V_j$ is contractible. The basis site has the same sheaf cohomology on
these basis objects as the corresponding topological subspace site by the
basis-site comparison in \Cref{lem:good-cover-site-tools}, so the constant
finite abelian sheaf $A$ has
\[
H^1(V_i,A)=0,
\qquad
H^1(V_i\cap V_j,A)=0.
\]
This is exactly the gerbe-adapted condition of
\Cref{def:gerbe-adapted-cech-comparison} for the constant band.

Finally, on a good cover all nonempty finite intersections are contractible.
Therefore the \v{C}ech cochains of the constant band $A$ are naturally the
simplicial cochains of the nerve $N(\mathcal V)$, with the same coboundary.
This gives $\check H^2(\mathcal V,A)\cong H^2(N(\mathcal V),A)$. For
$\mathcal U$, the nerve theorem recalled in
\Cref{lem:good-cover-site-tools} gives a homotopy equivalence
$|N(\mathcal U)|\simeq X$, hence the stated cohomology identification
\cite[Sec.~4G]{Hatcher2002}.
\end{proof}

\begin{corollary}[Concrete good-cover semantic models]\label{cor:good-cover-semantic-models}
Let $X$ be either $S^2$ or $\mathbb{R}P^2$, let $\mathcal U$ be a finite good cover of $X$, and let $A$ be a finite abelian group. Use the good-cover semantic slice $\mathcal C_X(\mathcal U)$ of \Cref{prop:good-cover-semantic-slice}, whose terminal object is $X$. For every class
\[
\omega\in H^2(N(\mathcal U),A)
\]
there is a one-reference semantic instance satisfying the gluing-sensitive
realization, branch \v{C}ech obstruction, and finite-nerve presentation
hypotheses such that:
\begin{enumerate}[label=(\roman*),nosep]
\item the chosen stack-banded realization prestack is the $A$-banded gerbe $\mathfrak G_\omega$;
\item the stackification morphism is the identity, hence global conservativity at the terminal fibre holds;
\item $\Vis_{p,s}(r)=\{v\}$ has one visible branch;
\item the branch obstruction class of $v$ is exactly $\omega$.
\end{enumerate}
If $\omega\neq 0$, then this instance satisfies
\[
M,p\Vdash \Null^{\mathrm{glue}}_s(r).
\]
In this case the example is genuinely nontrivial: the terminal fibre of the
chosen gerbe is empty by non-neutrality, while the gerbe is locally nonempty on
the good cover. Hence $F_{p,s}(a)=\varnothing$ but
\[
\bigl(a_{p,s}F_{p,s}\bigr)(a)\neq\varnothing,
\]
so compatible local sectionability and global sectionability are separated in
the semantic instance itself.
\end{corollary}

\begin{proof}
\Cref{prop:good-cover-semantic-slice} supplies the semantic slice with
terminal object $X$, finite common refinements, a cofinal filtered family of
gerbe-adapted good covers, and the local $H^1$-vanishing required by
\Cref{def:cech-gluing-obstruction}. It also identifies
$H^2(N(\mathcal U),A)$ with the \v{C}ech cohomology of the constant band on the
cover $\mathcal U$. Apply \Cref{prop:semantic-gerbe-realization} to the class
$\omega$. The resulting semantic instance uses the chosen gerbe
$\mathfrak G_\omega$ itself as the realization prestack, so the realization
hypothesis and global conservativity are verified by construction rather than
inferred from the band. If $\omega\neq 0$, the gerbe is non-neutral, and the
final assertion follows from \Cref{thm:gerbe-null-semantics}.
\end{proof}

\begin{example}[A visible quotient on an $S^2$-type nerve]\label{ex:s2-visible-quotient}
Let $\mathcal U$ be a finite good cover of $S^2$ with nerve $N(\mathcal U)$, and let
\[
A=\mathbb Z/4\mathbb Z.
\]
Finite intersections in a good cover are contractible when nonempty, so the cover is gerbe-adapted for the constant band $A$.
Since
\[
H_1(S^2,\mathbb Z)=0,
\qquad
H_2(S^2,\mathbb Z)\cong \mathbb Z
\]
\cite[Secs.~2.1, 2.2]{Hatcher2002}, the universal coefficient map is an isomorphism
\[
H^2(N(\mathcal U),A)\cong \mathrm{Hom}(H_2(N(\mathcal U),\mathbb Z),A)\cong A.
\]
Let $\omega$ be the class corresponding to the homomorphism sending a fundamental class to $2\in\mathbb Z/4\mathbb Z$. Then
\[
K_{\omega}=\langle 2\rangle=\{0,2\},
\qquad
A_{\mathrm{vis}}^{\omega}\cong (\mathbb Z/4\mathbb Z)/\langle 2\rangle\cong \mathbb Z/2\mathbb Z.
\]
Thus the obstruction is nontrivial, but only its parity quotient is intrinsically visible. By \Cref{cor:good-cover-semantic-models}, this class is realized by a chosen globally conservative stack-banded realization input with a single visible branch; in that realization \Cref{thm:gerbe-null-semantics} yields a genuine gluing-level absence with a nontrivial visible quotient.
\end{example}

\begin{example}[A character-blind obstruction on an $\mathbb{R}P^2$-type nerve]\label{ex:rp2-character-blind}
Let $\mathcal U$ be a finite good cover of $\mathbb{R}P^2$ with nerve $N(\mathcal U)$, and take
\[
A=\mathbb Z/2\mathbb Z.
\]
Again the good-cover condition gives the required local $H^1$-vanishing for the constant band $A$, while the nerve theorem identifies the homology of $N(\mathcal U)$ with that of $\mathbb{R}P^2$.
Then
\[
H_2(\mathbb{R}P^2,\mathbb Z)=0,
\qquad
H_1(\mathbb{R}P^2,\mathbb Z)\cong \mathbb Z/2\mathbb Z
\]
\cite[Sec.~2.2]{Hatcher2002}. The universal coefficient sequence therefore yields
\[
H^2(\mathbb{R}P^2,\mathbb Z/2\mathbb Z)
\cong
\operatorname{Ext}^1(\mathbb Z/2\mathbb Z,\mathbb Z/2\mathbb Z)
\cong
\mathbb Z/2\mathbb Z
\]
\cite[Sec.~3.1]{Hatcher2002}. Let $\omega$ be the nonzero class. Then
\[
\mathrm{ev}_{\omega}=0,
\qquad
K_{\omega}=0,
\qquad
A_{\mathrm{vis}}^{\omega}=A,
\]
by \Cref{thm:character-blind-obstructions}. Thus no coefficient-level quotient
collapse occurs: $A_{\mathrm{vis}}^\omega=A$ and every character of $A$ is
class-admissible. The blindness is instead at the obstruction-detection level:
because the evaluation map out of $H_2(\mathbb{R}P^2,\mathbb Z)$ is zero, all
character channels evaluate the nonzero class trivially. Hence this branch is
completely character-blind. By \Cref{cor:good-cover-semantic-models}, it occurs
as the unique visible branch of a chosen globally conservative stack-banded
realization input; for that semantic instance \Cref{thm:gerbe-null-semantics}
yields a genuine gluing-level absence that is invisible to every character
channel.
\end{example}


\subsection{Connection to sheaf-theoretic contextuality}

The unique-branch case $|\Vis_{p,s}(r)|=1$ deserves separate attention because it connects the present theory to the sheaf-theoretic analysis of contextuality and its cohomological refinements \cite{AbramskyBrandenburger2011,AbramskyMansfieldBarbosa2012,AbramskyBarbosaKishidaLalMansfield2015,Caru2017}. The comparison can be stated formally once the local object functor is specialized to the support presheaf of an Abramsky--Brandenburger measurement scenario.

\begin{theorem}[Unique-branch comparison with contextuality]\label{thm:unique-branch-contextuality-comparison}
Assume that $F_{p,s}$ is the support presheaf of an
Abramsky--Brandenburger scenario, equipped with a chosen gluing-sensitive
realization input satisfying the hypotheses of
\Cref{def:cech-gluing-obstruction} and global conservativity at $a$. Suppose
that
\[
\Vis_{p,s}(r)=\{v\}
\]
and that the unique branch $v$ admits a finite nerve presentation
\[
(\mathcal U,A,\omega).
\]
Then:
\begin{enumerate}[label=(\roman*),nosep]
\item compatible local sectionability of the support presheaf is equivalent to
the existence of the unique visible branch $v$;
\item global sectionability is equivalent to neutrality of the unique component
gerbe $\mathfrak G_v$, and gluing-level absence is equivalent to
non-neutrality of $\mathfrak G_v$;
\item if the chosen gerbe-adapted system $\mathcal A_v$ is $H^2$-detecting,
then
\[
M,p\Vdash \Null^{\mathrm{glue}}_s(r)
\quad\Longleftrightarrow\quad
\omega_v\neq 0;
\]
\item this no-global-section obstruction is invisible to every character
channel if and only if $\omega$ lies in the pure $\operatorname{Ext}$-summand
of the universal coefficient sequence for $H^2(N(\mathcal U),A)$.
\end{enumerate}
Thus the unique-branch support-presheaf setting recovers the
Abramsky--Brandenburger no-global-section obstruction. The Car\`u-type blind
cases are exactly the pure $\operatorname{Ext}$-residuals.
\end{theorem}

\begin{proof}
Because $F_{p,s}$ is the support presheaf of the measurement scenario,
its value on a context is the set of locally allowed outcome assignments, and
its restriction maps are the usual restriction of assignments to smaller
contexts. Under this specialization, a family of local sections over the
terminal context cover is matching in the sense of the support presheaf exactly
when it satisfies the common-refinement compatibility predicate
$\CompSec_s(r)$ used here; the existence of local supported assignments on the
cover is $\LocSec_s(r)$, and the existence of one supported assignment on the
terminal context is $\Sec_s(r)$. Thus compatible local sectionability means
exactly that the support presheaf has a compatible family over the terminal
context cover, while global sectionability means exactly that the support
presheaf has a global section. By
\Cref{thm:gerbe-null-semantics}, compatible local sectionability is equivalent
to $\Vis_{p,s}(r)\neq\varnothing$, so under the displayed uniqueness hypothesis
this is equivalent to the existence of the branch $v$. This proves
\textup{(i)}.

The same theorem identifies global sectionability with the existence of a
neutral visible component gerbe. Since there is only one visible branch, that
gerbe is $\mathfrak G_v$, and gluing-level absence is equivalently the
non-neutrality of $\mathfrak G_v$. This proves \textup{(ii)}.

If $\mathcal A_v$ is $H^2$-detecting, then
\Cref{prop:cech-comparison-detects-nonneutrality} identifies non-neutrality of
$\mathfrak G_v$ with nonvanishing of the branch class. This is \textup{(iii)}.

Finally, \Cref{thm:character-blind-obstructions} identifies complete
character-blindness of the finite nerve presentation $(\mathcal U,A,\omega)$
with the condition that $\omega$ lie in the image of the
$\operatorname{Ext}^1$-term. This gives \textup{(iv)} and the final sentence.
\end{proof}

\begin{remark}[Car\`u-type blind cases as pure $\operatorname{Ext}$-residuals]\label{rem:caru-ext-residual}
The pure-$\operatorname{Ext}$ characterization proved in
\Cref{thm:unique-branch-contextuality-comparison} turns the comparison with
cohomological contextuality methods
\cite{AbramskyMansfieldBarbosa2012,AbramskyBarbosaKishidaLalMansfield2015,Caru2017}
into a theorem: the failure of cohomological completeness in the unique-branch
case is not an external phenomenon but exactly the case in which the surviving
obstruction class lies entirely in the $\operatorname{Ext}$-summand of the
universal coefficient sequence.
\end{remark}

\begin{remark}[Conditional nature of the contextuality comparison]\label{rem:contextuality-comparison-conditional}
\Cref{thm:unique-branch-contextuality-comparison} is not a construction of a
stack-banded realization input from an arbitrary support presheaf. The support
presheaf supplies the set-valued local-object layer, while the gerbe, its
off-diagonal isomorphism sheaves, global conservativity, the stackification
banding, and the branch class are additional realization data required to apply
that cited theorem. Thus
the comparison with
Abramsky--Brandenburger contextuality and its cohomological refinements
\cite{AbramskyBrandenburger2011,AbramskyMansfieldBarbosa2012,AbramskyBarbosaKishidaLalMansfield2015}
is a conditional bridge: once such a realization has been supplied, the
no-global-section obstruction and the pure-$\operatorname{Ext}$ blind case are
identified by the displayed equivalences.
\end{remark}

